\section{Introduction}
%1 第一段 多智能体planning  价值 意义
%2 常用的几套方法，引用大论文
%3 nature science servery
%4 对比后 用ltl的优势，所以用ltl描述任务是很重要的部分，

Recent advances in computation, perception and communication
allow the deployment of autonomous robots in large, remote and hazardous environments,
e.g., to assist service staff in hospitals~\cite{jemal2015multi},
to maintain offshore drilling platforms~\cite{cliff2015online},
to monitor and assist construction sites~\cite{zhang2009collaborative}.
Furthermore, fleets of heterogeneous robots,
such as unmanned ground vehicles and unmanned aerial vehicles,
are deployed to accomplish tasks that are otherwise too inefficient
or even infeasible for a single robot~\cite{arai2002advances}.
Not only the overall efficiency of the team can be significantly improved
by allowing the robots to move and act concurrently~\cite{toth2002overview};
but also the capabilities of the team can be greatly extended by
enabling multiple robots to directly collaborate on a task~\cite{fink2008multi}.
Recent works have demonstrated such potentials preliminary for simple
tasks such as collaborative exploration~\cite{doi:10.1126/scirobotics.ade9548}, formation control~\cite{varava2017herding},
object transportation~\cite{doi:10.1126/scirobotics.abf1628} and pursuer-evader games \cite{doi:10.34133/research.0246}.
The task planning problem for multi-robot systems is in general NP-hard~\cite{1997Task},
due to the inherent combinatorial nature of robot-task assignment and various constraints such as capabilities and deadlines.
The standard approach is to formulate a mixed integer linear programs (MILP)
over the integer variables and constraints~\cite{agrawal2016scalable}.
Whereas being sound and optimal, these methods are applicable to only small-scale systems.
Thus, extensive work can be found on designing meta-heuristic algorithms for finding
sufficiently good solutions in a reasonable time,
e.g., genetic algorithms~\cite{KESHANCHI20171}, colony optimization~\cite{li2017multi,doi:10.34133/2020/1762107},
particle swarm optimization~\cite{yan2021task,8280872},
learning-based algorithms~\cite{wells2019learning,pmlr-v70-omidshafiei17a,liu2019task},
or large language models~\cite{}.\todo{add references}
However, these methods often lack a formal guarantee on the correctness and quality of the planning
results.


% Temporal tasks as LTL formulas
Moreover, to specify more complex tasks,
many recent works propose to use formal languages
such as Linear Temporal Logic (LTL)~\cite{Katoen2008Principles}, Computation Tree Logic (CTL)~\cite{koymans1990specifying},
and Signal Temporal Logic (STL)~\cite{maler2004monitoring}, as an intuitive yet
powerful way to describe both spatial and temporal requirements on
global~\cite{luo2021abstraction} or local~\cite{guo2016task} behaviors.
Notably, the works in~\cite{luo2021temporal, sahin2019multirobot, jones2019scratchs}
formulate MILP by a central planning unit given different system models and task constraints;
the works in~\cite{2016Decomposition,schillinger2018simultaneous,kantaros2020stylus}
instead propose various search algorithms over the state or solution space of the whole system.
However, the aforementioned planning methods are often executed {offline}
for a set of predefined {static} tasks.
A particularly challenging scenario is when the system operates indefinitely,
i.e., new tasks are released or canceled \emph{dynamically} and \emph{continually}
by external demand \cite{7306591};
or certain target features related to the tasks can change location
during run time~\cite{doi:10.34133/research.0246}.
This would require the fleet to adaptively change their task plans online to modify
existing assignments and incorporate new tasks.
Thus, the aforementioned methods become inadequate as the sequence of tasks is infinite
and their specifications are unknown beforehand.
Recursive application of the centralized methods in a naive way leads to
not only intractable computation complexity,
but also inconsistent or even oscillatory assignments.
Thus, an efficient and adaptive planning scheme is essential for multi-robot systems
that collaborate in a dynamic environment \cite{8901075,10102342,choudhury2022dynamic}
or an unknown environment \cite{9884983}.


%==============================
\subsection{Related Work} \label{related-work}
The standard framework for planning under temporal tasks is based on the model-checking algorithm~\cite{Katoen2008Principles}:
First, the task formulas are converted to a Deterministic Robin Automaton (DRA)
or Nondeterministic B\"uchi Automaton (NBA),
via off-the-shelf tools such as~\texttt{SPIN}~\cite{ben2010primer} and~\texttt{LTL2BA}~\cite{10.1007/3-540-44585-4_6}.
Second, a product automaton is created between the automaton of formula and the models of all agents,
such as weighted finite transition systems (wFTS)~\cite{Katoen2008Principles}, Markov decision processes~\cite{6702421}
or Petri nets~\cite{7294096}.
Last, certain graph search or optimization procedures are applied to find a valid and accepting plan
within the product automaton, such as nested-Dijkstra search~\cite{schillinger2018simultaneous}, integer programs~\cite{9663414},
auction~\cite{schillinger2019hierarchical} or sampling-based~\cite{kantaros2020stylus}.

Thus, the fundamental step of all aforementioned methods is to translate the task formula into
the associated automaton.
This translation may lead a double-exponential size
in terms of the formula length as shown in~\cite{10.1007/3-540-44585-4_6}.
The only exceptions are GR(1) formulas~\cite{10.1007/11609773_24}, of which the associated automaton
can be derived in polynomial time but only for limited cases.
%\todo{This is what I know. You can not say every type of LTL formula has double exponential complexity.}
In fact, for general LTL formulas with length more than~$25$,
it takes more than~$2$ hours and $13 GB$ memory to compute the associated NBA via~\texttt{LTL2BA}.
%目前最先进的几类方法和原始方法相比降低了很大的复杂度。但是他们的仿真算例使用的还是收到了NBA的限制。
Although recent methods have greatly reduced the planning complexity in other aspects,
the complexity of considered task specifications are still limited by this translation process.
%解释sampling 和 simultaneous， 规划的
For instance, the sampling-based method~\cite{kantaros2020stylus}
avoids creating the complete product automaton via sampling,
of which the largest simulated case has $400$ agents and
the task formula has a maximum length of $14$.
The planning method~\cite{schillinger2018simultaneous} decomposes the resulting task automaton
into independent sub-tasks for parallel execution.
The simulated case scales up to~$100$ robots and a task formula of length~$18$.
% 其他算法的平均表现
Moreover, other existing works such
as~\cite{2021Reactive,2018Multi,6942756,7487762} mostly
{consider} task formulas of length around~$6-10$.
% 因此，这个限制了在物理世界的应用
This limitation hinders greatly its application to more complex robotic missions.
%where the planning process has the time budget of several seconds.
%\todo{I think you can add more detailed comprisons with other work,  as this part is quite short.}


%==============================

This drawback becomes even more apparent for dynamic scenes,
where contingent tasks specified as LTL formulas are triggered by external observations
and released to the whole team online.
In such cases, most existing approaches compute
the automaton associated with the new task, of which the synchronized
{\emph{product} with the current automaton is derived, see e.g.,~\cite{6942756, 7487762}}.
Thus, the size of the resulting automaton equals to the product
of \emph{all} automata associated with the contingent tasks,
which is clearly a combinatorial blow-up.
%真找不到了，所有相关的算法只有15,16.人家也不说自己的复杂度爆炸的事情
Consequently, the amount of contingent tasks that can be
handled by the aforementioned
approaches is limited to hand-picked examples.

\subsection{Our Contribution}
To overcome this curse of dimensionality in the size of tasks and agents,
we propose a new paradigm that is \emph{fundamentally} different from
the model-checking-based methods.
First, for a LTL formula that consists of numerous sub-formulas of smaller size,
we calculate the R-posets of each sub-formula as a set of subtasks and
their partial temporal constraints.
Then, an efficient algorithm is proposed to compute the product of
R-posets associated with each sub-formula.
{The resulting product of R-posets is complete in the sense}
that it retains all subtasks from each R-poset along with their partial orderings
and resolves potential conflicts.
Given this product,
a task assignment algorithm called the time-bound contract net (TBCN) is proposed to
assign collaborative subtasks to the agents,
subject to the partial ordering constraints.
Last but not least, the same algorithm is applied online to dynamic scenes
where contingent tasks are
triggered and released by online observations.
It is shown formally that the proposed method has a polynomial time
and memory complexity to derive the first valid plan
w.r.t. the system size and formula length.
Extensive large-scale simulations and experiments are conducted
for a hospital environment where service robots react
to online requests such as collaborative transportation,
cleaning and monitoring.

%具体介绍方法，有哪些贡献 特点 更技术一点


%上段技术细节
%下端单纯讲效果  效果1 能处理大系统 2 考虑interactive objects 3 能反应环境特征
Main contribution of this work is three-fold:
(i) a systematic method is proposed to tackle task formulas with length more than~$400$,
which overcomes the limitation of existing translation tools that can only process
formulas of length less than~$25$ in reasonable time;
(ii) an efficient algorithm is proposed to decompose and integrate
contingent tasks that are released online,
which not only avoids a complete re-computation of the task automaton
but also ensures a polynomial complexity to derive the first valid plan;
(iii) the proposed task assignment algorithm, TBCN, is fully compatible with
both static and dynamic scenarios with interactive objects.

%==============================
\subsection{Problem Statement}
{Consider a multi-agent system with heterogeneous capabilities,
  a series of LTL formulas that are released online,
  and a set of interactive objects in the dynamic environment.
  The objective is to generate a trajectory plan for the system such that
  these tasks are satisfied with high efficiency.}


{For instance as shown in Sec.~\ref{simulation},
  a fleet of heterogeneous service robots is deployed in the hospital environment.
  Different tasks such as transportation of goods or patients, cleaning and maintenance
  are released online continuously.
  Many of such tasks contain numerous subtasks  with ordering constraints
  that require direct collaboration of different robots.
  An efficient coordination algorithm is proposed such that these subtasks
  are assigned and fulfilled online in a timely manner.
}



%% Our approach is based on a temporal logic, called Linear Temporal Logic(LTL), which provides an
%% accurate, expressive and easily understandable way to specify and tasks (e.g. "take the vomit illness into
%% operating room) and safety requirement ("do not enter vomiting area until it is disinfected"),
%% and the potential prior knowledge ("facts that the operation cannot be executed before the patient is taken
%% to the operating room."). We converted each formula into an R-poset, which contains the sutbasks and the
%% temporal order between them, and proposed a Product method between R-posets that can get a sub-optimal solution
%%  in linear complexity. Then, the action sequences are generated by operation research algorithm to satisfy both
%%  tasks and objects requirement.
%%  \todo{I do not see the purpose of this paragraph.
%%    Isn't the description of method already introduced in the second last paragraph of the introduction?
%%  I would have to repeat the same thing here?}
